\maketitle

\begin{abstract}
A system that improves the machinery it uses to acquire capabilities raises two distinct questions: whether the modification helps, and where it stops helping. We report six frozen experiments, M107--M112, from the public \genesis research program. The tested lineage and qualification arms are model-free; M112 uses one separately isolated model invocation only to generate a sealed held-out world bank. M107 acquires an operator that expands the lower interpreter's complete image from 4 to 16 Boolean functions. M108 then acquires a state-held attribution rule that changes the machinery performing later acquisitions. M109 replaces authored blame labels with controlled trials and shows two successive machinery generations, with constructive reach increasing $6\subset20\subset243$. Restored unchanged into a materially different carrier, the acquired cascade transfers positively on the tested rows inside the producer's reachable failure census. On a realizable row outside that census it becomes confidently wrong: on six of six row-5 demands a fresh predecessor succeeds while both descendants fail, even though \reach grows strictly on every world. M111 recovers part of this loss by acquiring, at a third machinery generation, a diagnostic policy that spends one scarce probe only where its feature vocabulary does not determine an answer; never-probe, always-probe, and static strategies all fail one member of the ambiguous pair under the same budget. Finally, M112 removes direct project selection and authorship of world contents through a single blind, sealed generator invocation. Diagnosis reproduces 24/24 and every scientific transfer outcome reproduces, including the harm, while the inherited transfer checker remains negative at 22/24 because one blind world violates a fixed-point certificate. No generality gate advances. The contribution is bounded, auditable evidence of acquisition-machinery adaptation, census-conditional negative transfer, and self-directed diagnosis, not open-ended recursive self-improvement, training-data independence, independent reproduction, or AGI.
\end{abstract}

\section{Introduction}
Recent self-improving agent systems increasingly modify not only task solutions but also reusable skills, agent programs, prompts, memory structures, or meta-level improvement procedures \cite{yin2024godelagent,robeyns2025selfimproving,zhang2025dgm,zhang2026hyperagents,wang2026metaskill,ren2026survey}. This direction inherits an old aspiration from formally self-referential systems such as the G{\"o}del Machine \cite{schmidhuber2007godel}, but modern systems usually substitute empirical evaluation for proof and foundation-model competence for a closed symbolic substrate. Their practical successes make a narrower measurement problem increasingly important: \emph{what exactly improved, what later capability causally depends on that improvement, and under what conditions can the improvement become harmful?}

Transfer learning has long recognized that source-derived structure can improve a target task or make it worse when source and target differ in a relevant way \cite{rosenstein2005transfer,pan2010transfer,wang2019negative,zhang2020negative}. The problem is sharper for self-modifying systems because a change to an acquisition procedure is not merely one more task-specific skill. It becomes an inductive bias applied to later learning and adaptation. A procedure that was locally beneficial can therefore enlarge the space of things a lineage can construct while simultaneously making its decisions worse outside the failure geometry that produced it.

This paper studies that possibility in \genesis, a public, append-only experimental program on persistent adaptive software. Genesis I reported M094--M100, culminating in bounded causal cumulative capability acquisition through serialized state and hard producer-process death \cite{mets2026genesis1}. The present paper begins where that work stopped. Instead of asking whether an acquired task operation can enable a later operation, we ask whether the lineage can change the \emph{machinery that decides how later capabilities are acquired}, whether that changed machinery can produce further machinery changes, whether those changes transfer to a materially different carrier, and whether the lineage can detect when its inherited decision vocabulary does not determine the right action.

The central result is deliberately not a success. M110 restores a two-generation acquisition-machinery cascade unchanged into a carrier that took no part in producing it. Inside the producer's reachable attribution census the cascade adds capability. Outside it, the same acquired rule fires confidently on a row the producer could never present and points to the wrong component. A fresh predecessor solves all six such demands; both descendants solve none. Yet the descendants' constructive reach grows monotonically on every world. Figure~\ref{fig:capacity} summarizes the dissociation: capacity rises while realized competence does not.

M111 then asks whether the lineage can use its own historical record to identify an underdetermined region and spend a scarce experiment there. It acquires one diagnostic policy from a pooled record. The policy does not invent experimentation; the probe primitive and budget remain authored. It does, however, preserve its single probe on a determined demand and spend it on an ambiguous one, succeeding under a budget that an always-probe policy wastes. The policy is inexpressible one machinery generation earlier by a monotonicity lemma, so the previous generation creates the language in which the next diagnostic modification can exist.

